\section*{Ethical Considerations}
Local processing can reduce document disclosure, but it does not eliminate privacy risk: PDFs, extracted text, figures, embeddings, traces, prompts, and model outputs remain on the device and require deletion, access-control, and retention practices appropriate to the user's environment. The interface discloses the active network mode, and the paper distinguishes external acquisition from strict-local analysis. Benchmark documents and model artifacts must be used under their licenses. Human claim annotation and the researcher pilot must not begin before the responsible institution provides the applicable ethics or IRB determination; consent, compensation, de-identification, exclusions, and data retention must then follow the approved process. Citation-like output may create unwarranted trust, so the system exposes source evidence and selective failure and must not market a retrieved page as proof that a statement is true.
